% ============================================================================
% Fichier  : partie5/P05_C19_gouvernance_donnees_ia.tex
% Titre    : Gouvernance des Données et de l'Intelligence Artificielle
% Auteur   : MINKA MI NGUIDJOÏ Thierry Emmanuel
% Version  : 1.0 -- octobre 2026
% ============================================================================
% Sources : COBIT 2019, fiche APO14 (dossier Normes) ; ISO/IEC 38507:2022,
% ISO/IEC 42001:2023, ISO/IEC 23894:2023 (présentation d'après leur objet
% public) ; NIST AI RMF 1.0 (2023) ; Recommandation de l'UNESCO sur
% l'éthique de l'IA (2021) ; règlement (UE) 2024/1689. Loi n° 2024/017 :
% aucun article cité ; le régime des décisions automatisées est à vérifier
% dans le texte publié.
% ============================================================================

\chapter{Gouvernance des Données et de l'Intelligence Artificielle}
\label{chap:donnees-ia}
\chaptermark{Données et intelligence artificielle}

\epigraph{%
    « Un algorithme ne décide jamais seul.\\
    Quelqu'un a décidé de le laisser décider. »%
}{Principe de la gouvernance de l'IA}

\niveauChapitre{Expert}{%
    Maîtrise des Chapitres~\ref{chap:droit} (protection des données
    personnelles) et~\ref{chap:agile-cloud} (informatique en nuage).%
}

\begin{boxobjectifs}
\begin{itemize}
    \item Définir la gouvernance des données (propriété, qualité, cycle de
          vie).
    \item Présenter l'ISO/IEC~38507 et l'ISO/IEC~42001.
    \item Évaluer les risques d'un système d'IA (biais, explicabilité,
          responsabilité).
    \item Appliquer la loi de 2024 sur les données personnelles à un cas de
          notation par IA.
\end{itemize}
\end{boxobjectifs}

\bigskip

% ============================================================================
\section*{Question directrice : qui répond d'une décision prise par un modèle~?}
% ============================================================================

Le Chapitre~\ref{chap:cas-integrateurs} l'a établi~: le modèle de notation
de crédit de NOUN PAY refuse 61~\% des demandes d'une région, contre
24~\% en moyenne nationale. Il utilise l'opérateur mobile du demandeur,
très corrélé à la région. Personne, dans l'entreprise, n'était chargé
d'évaluer les effets du modèle sur les clients (Chapitre~\ref{chap:processus}).
L'Autorité de protection des données demande des comptes.

Ce cas réunit les deux sujets de ce chapitre. La \textbf{gouvernance des
données} (sections~\ref{sec:c19-donnees} et~\ref{sec:c19-qualite}) traite
de ce que l'organisation possède~: qui en répond, avec quelle qualité, pour
combien de temps. La \textbf{gouvernance de l'intelligence artificielle}
(sections~\ref{sec:c19-ia} à~\ref{sec:c19-noun}) traite de ce que
l'organisation en fait quand elle confie des décisions à des modèles.

\separateur

% ============================================================================
\section{La Gouvernance des Données}
\label{sec:c19-donnees}
% ============================================================================

\subsection{Les données, un actif}

\begin{boxdef}
\termecle{Gouvernance des données}\index{gouvernance des données}~:
exercice de l'autorité et du contrôle sur la gestion des données~:
qui décide de leur définition, de leur usage, de leur qualité, de leur
protection et de leur conservation, selon quelles règles, et qui en rend
compte.
\end{boxdef}

COBIT~2019 consacre aux données un objectif de management à part entière,
APO14, qui couvre tout le cycle de vie des données, de leur création à
leur archivage. Ses liens principaux sont la qualité de l'information
financière (EG04) et de l'information de gestion (EG07, AG10)~: les
données sont d'abord ce sur quoi l'on décide.

\begin{boxobjet}{APO14 · Gérer les données}
\textbf{Description~:} gérer efficacement les actifs de données de
l'organisation sur tout leur cycle de vie, de la création à la diffusion,
à la maintenance et à l'archivage.

\textbf{Finalité~:} assurer l'utilisation efficace des données critiques
pour atteindre les objectifs de l'organisation.

\textbf{Objectifs liés~:} EG04, EG07~; AG10.

\textbf{Pratiques~:} APO14.01 définir la stratégie de gestion des données
et les rôles~; APO14.02 tenir un glossaire métier~; APO14.03 gérer les
métadonnées~; APO14.04 définir une stratégie de qualité des données~;
APO14.05 profiler les données~; APO14.06 évaluer leur qualité~; APO14.07
définir leur nettoyage~; APO14.08 gérer le cycle de vie des actifs de
données~; APO14.09 organiser l'archivage et la conservation~; APO14.10
gérer les sauvegardes et les restaurations.

\textbf{Indicateurs~:} part des données critiques dotées d'un
propriétaire, taux d'anomalies de qualité, part des données conservées
au-delà de leur durée prévue.
\end{boxobjet}
{\small\itshape Description, finalité, objectifs liés et pratiques résumés
et traduits par nos soins à partir de COBIT~2019~\cite{isaca2018cobitobjectifs}.}

\subsection{Les rôles}

\begin{tableaucours}{L{3.4cm}L{5.0cm}L{3.8cm}}{%
    \tableauHeader{Rôle} &
    \tableauHeader{Responsabilité} &
    \tableauHeader{Chez NOUN PAY}
}
Propriétaire des données & Répond d'un domaine de données~: usage
autorisé, niveau de qualité, accès & Directeur commercial pour les données
clients \\
Intendant des données (\textit{data steward}) & Veille au quotidien à la
définition et à la qualité & Analyste du crédit \\
Dépositaire & Assure le stockage, la sécurité et la sauvegarde &
Directeur technique et fournisseur de nuage \\
Délégué à la protection des données & Veille au respect du droit des
données personnelles & M.~DJOUMESSI \\
\end{tableaucours}
\vspace{-0.8em}
\begin{center}{\small\itshape Tableau~19.1. Les rôles de la gouvernance des
données}\end{center}

La règle cardinale~: \textbf{le propriétaire des données est un métier,
pas l'informatique}. La DSI conserve et protège~; c'est le métier qui
décide de ce que les données signifient et de l'usage qui en est fait.

% ============================================================================
\section{Qualité et Cycle de Vie}
\label{sec:c19-qualite}
% ============================================================================

\subsection{Les dimensions de la qualité}

\begin{tableaucours}{L{3.0cm}L{5.0cm}L{4.2cm}}{%
    \tableauHeader{Dimension} &
    \tableauHeader{Question} &
    \tableauHeader{Exemple}
}
Exactitude & La donnée reflète-t-elle la réalité~? & Le solde converti
est-il le bon (SANAGA)~? \\
Complétude & Toutes les données requises sont-elles présentes~? &
Revenus manquants dans 18~\% des dossiers de crédit \\
Cohérence & La même donnée a-t-elle la même valeur partout~? & L'adresse
du client dans deux systèmes \\
Actualité & La donnée est-elle à jour~? & Position d'un conteneur vieille
de six heures (KRIBILOG) \\
Validité & Respecte-t-elle le format et les règles~? & Numéro de
téléphone au bon format \\
Unicité & Chaque entité n'apparaît-elle qu'une fois~? & Un fournisseur
enregistré trois fois à l'ANPP \\
\end{tableaucours}
\vspace{-0.8em}
\begin{center}{\small\itshape Tableau~19.2. Les dimensions de la qualité
des données}\end{center}

\subsection{Le cycle de vie}

Une donnée se crée, s'utilise, se partage, s'archive et se détruit. La
gouvernance fixe, pour chaque catégorie, sa \emph{durée de conservation}~:
assez longue pour répondre aux obligations (dix ans pour les données de
connexion selon l'article~29 de la loi n°~2010/012, Chapitre~\ref{chap:droit}),
pas plus longue que nécessaire pour les données personnelles. Conserver
« au cas où » n'est pas une politique, c'est un risque.

\begin{boxprocedure}
\textbf{Le registre des données critiques.} Pour qu'un auditeur puisse
établir que les données critiques sont gouvernées, l'organisation tient
un registre qui indique, pour chaque domaine de données~:
\begin{itemize}
    \item le propriétaire métier et l'intendant~;
    \item la définition, dans le glossaire métier~;
    \item la classification (publique, interne, confidentielle,
          personnelle, sensible)~;
    \item la localisation (systèmes, pays)~;
    \item les règles de qualité et leur dernier contrôle~;
    \item la durée de conservation et son fondement.
\end{itemize}
\end{boxprocedure}

% ============================================================================
\section{Gouverner l'Intelligence Artificielle}
\label{sec:c19-ia}
% ============================================================================

\subsection{Ce que l'IA change}

Un système d'intelligence artificielle n'est pas un logiciel comme un
autre, pour trois raisons qui intéressent directement la gouvernance~:
\begin{itemize}
    \item \textbf{il apprend des données}~: ses erreurs viennent souvent de
          biais présents dans les données d'entraînement, que personne n'a
          voulus~;
    \item \textbf{il est difficile à expliquer}~: certains modèles
          produisent un résultat sans que l'on puisse dire simplement
          pourquoi~;
    \item \textbf{il évolue}~: un modèle réentraîné est un nouveau modèle,
          et ses effets peuvent dériver avec le temps.
\end{itemize}

\subsection{Les normes et cadres de référence}

\begin{tableaucours}{L{3.4cm}L{2.0cm}L{6.8cm}}{%
    \tableauHeader{Référence} &
    \tableauHeader{Nature} &
    \tableauHeader{Objet}
}
ISO/IEC 38507 (2022) & Lignes directrices & Implications, pour l'organe de
gouvernance, de l'usage de l'IA par l'organisation~; prolonge
l'ISO/IEC~38500 \\
ISO/IEC 42001 (2023) & Norme certifiable & Système de management de l'IA,
selon la structure harmonisée (Chapitre~\ref{chap:normes}) \\
ISO/IEC 23894 (2023) & Lignes directrices & Management des risques liés à
l'IA, dans la logique de l'ISO~31000 \\
NIST AI RMF 1.0 (2023) & Cadre & Quatre fonctions~: gouverner,
cartographier, mesurer, gérer~\cite{nist2023airmf} \\
Recommandation de l'UNESCO (2021) & Instrument normatif international &
Valeurs et principes éthiques de l'IA, adoptés par les États
membres~\cite{unesco2021ethique} \\
\end{tableaucours}
\vspace{-0.8em}
\begin{center}{\small\itshape Tableau~19.3. Normes et cadres de la
gouvernance de l'IA}\end{center}

La Recommandation de l'UNESCO sur l'éthique de l'intelligence
artificielle, adoptée en novembre~2021, est le premier instrument normatif
mondial sur le sujet. Elle insiste sur la proportionnalité, la sécurité,
l'équité et la non-discrimination, la transparence et l'explicabilité, la
responsabilité et la surveillance humaine. Pour une formation dispensée par
une Chaire UNESCO, elle est le point de départ naturel.

\begin{boxinfo}
\textbf{Le contexte africain et international.} L'Union africaine a adopté
en 2024 une stratégie continentale sur l'intelligence artificielle.
L'Union européenne a adopté le règlement (UE)~2024/1689 sur l'intelligence
artificielle, qui classe l'évaluation de la solvabilité des personnes
physiques parmi les systèmes à haut risque~\cite{ue2024aiact}. Ce règlement
ne s'applique pas au Cameroun en tant que tel, mais il s'impose aux
entreprises qui y servent des clients, et sert déjà de référence aux
investisseurs.
\end{boxinfo}

% ============================================================================
\section{Évaluer les Risques d'un Système d'IA}
\label{sec:c19-risques}
% ============================================================================

\subsection{Une grille en six risques}

\begin{tableaucours}{L{3.0cm}L{5.0cm}L{4.2cm}}{%
    \tableauHeader{Risque} &
    \tableauHeader{Question} &
    \tableauHeader{Mesure de maîtrise}
}
Biais et discrimination & Le modèle traite-t-il défavorablement un groupe~?
& Mesure des écarts entre groupes, revue des variables \\
Explicabilité & Peut-on expliquer une décision à la personne concernée~?
& Motifs principaux de chaque décision, documentation du modèle \\
Responsabilité & Qui répond d'une décision du modèle~? & Un A dans la
matrice RACI (Chapitre~\ref{chap:processus}) \\
Robustesse & Le modèle résiste-t-il aux données inhabituelles et aux
manipulations~? & Tests de robustesse, surveillance des entrées \\
Dérive & Ses performances se dégradent-elles avec le temps~? & Indicateurs
suivis en production, réentraînement contrôlé \\
Conformité & Les données sont-elles traitées licitement~? & Analyse au
regard de la loi n°~2024/017 \\
\end{tableaucours}
\vspace{-0.8em}
\begin{center}{\small\itshape Tableau~19.4. Les risques d'un système d'IA}\end{center}

\subsection{Mesurer un biais}

Une mesure simple de l'écart entre groupes est le \termecle{ratio d'impact
disproportionné}\index{ratio d'impact disproportionné}~: le taux de
décisions favorables du groupe le moins favorisé, divisé par celui du
groupe le plus favorisé. Une pratique répandue, issue du droit américain
du travail, considère comme suspect un ratio inférieur à~0,8 (la « règle
des quatre cinquièmes »). Ce seuil n'a pas de valeur juridique au
Cameroun~; il sert de signal d'alerte.

\begin{boxscenario}{Opération NOUN. Le calcul du biais}
\textbf{Données.} Taux de refus de 61~\% dans la région concernée, soit un
taux d'acceptation de 39~\%. Taux de refus moyen national de 24~\%, soit un
taux d'acceptation de 76~\%.

\medskip
\textbf{Calcul.} Ratio d'impact $= 0{,}39 / 0{,}76 \approx 0{,}51$, très en
dessous du seuil d'alerte de 0,8.

\medskip
\textbf{Constatation.} Le ratio est de 0,51~; la variable « opérateur
mobile », corrélée à la région, figure parmi les variables du modèle.

\medskip
\textbf{Analyse} (confiance moyenne). L'écart peut s'expliquer en partie
par des différences réelles de risque, mais l'usage d'une variable
corrélée à la région fait peser un risque sérieux de discrimination
indirecte. Seule une analyse du modèle sans cette variable, à risque égal,
permettra de conclure. En attendant, la variable doit être suspendue
(Chapitre~\ref{chap:cas-integrateurs}).
\end{boxscenario}

% ============================================================================
\section{La Notation de NOUN PAY au Regard du Droit}
\label{sec:c19-noun}
% ============================================================================

\subsection{Les questions à instruire}

La loi n°~2024/017 (Chapitre~\ref{chap:droit}) s'applique aux données
personnelles utilisées par le modèle. Sans préjuger de la lettre de ses
articles, que l'on vérifiera dans le texte publié, six questions doivent
être instruites~:
\begin{enumerate}
    \item \textbf{La base juridique}~: sur quel fondement le traitement des
          données des demandeurs est-il licite~?
    \item \textbf{La finalité et la proportionnalité}~: les données utilisées
          sont-elles nécessaires à l'évaluation du crédit~? L'opérateur
          mobile l'est-il~?
    \item \textbf{L'information des personnes}~: les demandeurs savent-ils
          que leur demande est évaluée par un modèle~?
    \item \textbf{Les décisions automatisées}~: de nombreuses lois de
          protection des données de cette génération encadrent les
          décisions fondées exclusivement sur un traitement automatisé et
          ouvrent un droit à l'intervention humaine~; il faut vérifier si
          et comment la loi camerounaise le prévoit.
    \item \textbf{Le transfert hors du territoire}~: les données sont
          hébergées hors du Cameroun (Chapitre~\ref{chap:agile-cloud}).
    \item \textbf{La sécurité et la conservation}~: qui accède aux données,
          et combien de temps sont-elles gardées~?
\end{enumerate}

\begin{boxalert}
\textbf{Violation typique.} Une entreprise supprime de son modèle la
variable « région » pour éviter toute discrimination, mais conserve des
variables qui lui sont étroitement corrélées (opérateur mobile, préfixe
téléphonique, agence). Le biais demeure, par une autre voie.

\smallskip
\textbf{Règle~:} on ne prouve pas l'absence de biais en retirant une
variable~; on la prouve en \emph{mesurant} les résultats par groupe.
\end{boxalert}

\subsection{Le dispositif de gouvernance}

\begin{boxmetier}{Responsable GRC}
\textbf{Doter NOUN PAY d'une gouvernance de l'IA proportionnée}
\begin{enumerate}
    \item \textbf{Un inventaire}~: lister les modèles en production, leur
          usage, les données qu'ils emploient et leurs effets sur les
          clients.
    \item \textbf{Une classification}~: la notation de crédit, qui décide
          de l'accès d'une personne à un service financier, est à risque
          élevé.
    \item \textbf{Un responsable}~: désigner un A pour chaque modèle à
          risque élevé, au niveau de la direction (Chapitre~\ref{chap:processus}).
    \item \textbf{Une évaluation d'impact avant chaque version}~: mesure
          des écarts entre groupes, avis du délégué à la protection des
          données, décision écrite.
    \item \textbf{Une surveillance en production}~: ratio d'impact et taux
          d'erreur suivis chaque mois, avec un seuil qui déclenche la
          suspension.
    \item \textbf{Un recours humain}~: toute personne refusée peut demander
          un réexamen par un analyste.
\end{enumerate}
\end{boxmetier}

\begin{boxmetier}{Auditeur SI}
\textbf{Auditer un modèle de décision}
\begin{enumerate}
    \item Demander la documentation du modèle~: finalité, données
          d'entraînement, variables, performances, limites connues.
    \item Vérifier que chaque mise à jour a fait l'objet d'une évaluation
          d'impact et d'une autorisation (BAI06).
    \item Recalculer, sur un échantillon de décisions récentes, les taux
          par groupe et le ratio d'impact.
    \item Vérifier l'existence et l'usage effectif du recours humain.
\end{enumerate}
\end{boxmetier}

\begin{boiteRemarque}{L'humain dans la boucle suffit-il~?}
Prévoir qu'un analyste « valide » chaque décision du modèle est souvent
présenté comme la garantie ultime. L'expérience montre pourtant qu'un
analyste qui traite deux cents dossiers par jour finit par entériner les
propositions du modèle. La surveillance humaine n'est réelle que si
l'analyste dispose du temps, de l'information (les motifs de la
proposition) et du pouvoir de s'écarter du modèle sans avoir à s'en
justifier davantage que s'il le suivait.
\end{boiteRemarque}

% ============================================================================
\separateur
% ============================================================================

\begin{boxresume}
\textbf{Ce qu'il faut retenir de ce chapitre~:}
\begin{itemize}
    \item \textbf{Les données sont un actif gouverné}~: un propriétaire
          métier, un intendant, un dépositaire~; COBIT y consacre
          l'objectif APO14.
    \item \textbf{Six dimensions de qualité} (exactitude, complétude,
          cohérence, actualité, validité, unicité) et un cycle de vie avec
          des durées de conservation fondées.
    \item \textbf{L'IA apprend, s'explique mal et dérive}~: elle appelle une
          gouvernance propre.
    \item \textbf{Les références}~: ISO/IEC~38507 pour le conseil,
          ISO/IEC~42001 pour le système de management, ISO/IEC~23894 pour
          les risques, le NIST AI RMF, et la Recommandation de l'UNESCO
          pour les principes éthiques.
    \item \textbf{Mesurer le biais}~: le ratio d'impact de NOUN PAY (0,51)
          est un signal d'alerte~; retirer une variable ne prouve rien.
    \item \textbf{Un modèle à risque élevé} a un responsable, une évaluation
          d'impact à chaque version, une surveillance en production et un
          recours humain réel.
\end{itemize}
\end{boxresume}

\bigskip

\begin{boxexo}[title={\ding{45}\quad Exercice~19.1. Quelle dimension de qualité~? \niveaubase}]
Rattachez chacune des anomalies suivantes à une dimension de qualité
(Tableau~19.2).
\begin{enumerate}[(a)]
    \item Un agent payé deux fois par l'ANPP parce qu'il figure sous deux
          matricules.
    \item Un client de la BANQUE SANAGA dont la date de naissance est le
          31~février.
    \item Un conteneur localisé à Kribi alors qu'il est à Douala depuis la
          veille.
    \item Un dossier de crédit sans les revenus du demandeur.
\end{enumerate}
\end{boxexo}

\begin{boxexo}[title={\ding{45}\quad Exercice~19.2. Le registre des données de la BANQUE SANAGA \niveaumoyen}]
Établissez les cinq premières lignes du registre des données critiques de
la BANQUE SANAGA (exigence de contrôle de la section~\ref{sec:c19-qualite}),
en commençant par les soldes des comptes.
\begin{enumerate}[(a)]
    \item Pour chaque domaine, désignez le propriétaire métier et
          justifiez.
    \item Fixez les règles de qualité et leur fréquence de contrôle.
    \item Indiquez la durée de conservation et son fondement, en signalant
          ce qui reste à vérifier.
\end{enumerate}
\end{boxexo}

\begin{boxexo}[title={\ding{45}\quad Exercice~19.3. La réponse à l'Autorité de protection des données \niveauexpert}]
Rédigez, pour Mme~KAMGA, la réponse de NOUN PAY à l'Autorité de
protection des données.
\begin{enumerate}[(a)]
    \item Exposez les faits, en séparant constatations et analyses.
    \item Répondez aux six questions de la section~\ref{sec:c19-noun}, en
          indiquant ce qui suppose une vérification du texte de la loi.
    \item Présentez les mesures prises et prévues, avec leur échéance.
    \item Expliquez en quoi la Recommandation de l'UNESCO sur l'éthique de
          l'IA éclaire les engagements pris.
\end{enumerate}
\end{boxexo}

\bigskip

\begin{boxinfo}
\textbf{Transition.} La Partie~\ref{part:avance} a confronté la gouvernance
des SI aux défis de la vitesse, du nuage, des données et de
l'intelligence artificielle. La Partie~\ref{part:recherche} prend du
recul~: le Chapitre~\ref{chap:recherche-mesure} présente l'état de la
recherche sur la mesure de la gouvernance des SI.
\end{boxinfo}
